\chapter{เคมีสารสนเทศและการค้นหายาใหม่ด้วย ChEMBL, PubChem และ PyMOL}

กระบวนการค้นหายาและพัฒนาสารออกฤทธิ์ทางชีวภาพ (Drug Discovery) ได้รับประโยชน์อย่างมหาศาลจากการบูรณาการโมเดลเคมีเชิงคำนวณเข้ากับฐานข้อมูลชีวโมเลกุลขนาดใหญ่

\section{ฐานข้อมูลสารออกฤทธิ์ทางชีวภาพ ChEMBL}

ChEMBL รวบรวมข้อมูลโมเลกุลขนาดเล็กที่มีการทดลองวัดค่าการยับยั้งเอนไซม์หรือตัวรับสัญญาณ (Bioactivity Assays) โดยมีตัวชี้วัดสำคัญ ได้แก่:
\begin{itemize}
    \item \textbf{IC50 (Half Maximal Inhibitory Concentration):} ความเข้มข้นของสารที่ยับยั้งการทำงานของโปรตีนได้ร้อยละ 50
    \item \textbf{Ki (Inhibition Constant):} ค่าคงที่การยับยั้งทางอุณหพลศาสตร์
    \item \textbf{Lipinski's Rule of Five:} กฎ 5 ข้อสำหรับประเมินความเป็นไปได้ในการพัฒนาเป็นยารับประทาน
\end{itemize}

\section{การสืบค้นโครงสร้างและคุณสมบัติทางเคมีด้วย PubChem}

PubChem เป็นคลังข้อมูลสารเคมีที่ใหญ่ที่สุดในโลก การสืบค้นสามารถทำผ่านรหัสระบุชนิดสาร (CID), รหัสโครงสร้างแบบบรรทัดเดียว (SMILES), หรือการเปรียบเทียบความคล้ายคลึงของโครงสร้างโมเลกุล (Tanimoto Similarity Coefficient):

\begin{equation}
T(A, B) = \frac{|A \cap B|}{|A \cup B|}
\end{equation}

\section{การสร้างภาพจำลองปฏิสัมพันธ์โมเลกุลด้วย PyMOL}

เมื่อได้สารที่มีศักยภาพแล้ว ซอฟต์แวร์ PyMOL จะถูกนำมาใช้สร้างแบบจำลอง 3 มิติเพื่อแสดงพันธะไฮโดรเจน (Hydrogen Bonds), แรงดึงดูดทางไฟฟ้าสถิต (Electrostatic Interactions), และโพรงที่จับของยา (Binding Pocket Surface) ช่วยให้นักวิจัยเข้าใจกลไกการออกฤทธิ์ในระดับชีวฟิสิกส์ได้อย่างลึกซึ้ง
